\subsection{2. Những nghiên cứu liên quan}
\indent\indent Những năm gần đây, việc khai thác dữ liệu đa phương thức đang trở thành xu hướng chủ đạo. Sự kết hợp giữa thông tin đặc trưng từ hình ảnh, âm thanh và ngữ nghĩa văn bản giúp máy tính "hiểu" dữ liệu ICH một cách trọn vẹn hơn. Tuy nhiên, thách thức lớn nhất vẫn nằm ở việc mô hình hóa sự tương tác phức tạp giữa các phương thức dữ liệu khác nhau.\vspace{0.3cm}

Trong lĩnh vực phân loại ICH, Qin và cộng sự \cite{FAN2023120555} đã chỉ ra rằng các phương pháp truyền thống thường chỉ tập trung vào đặc trưng thị giác mà bỏ qua các mô tả văn bản đi kèm. Để khắc phục, nhóm tác giả đề xuất mô hình MICMLF (Multimodal Attention and Hierarchical Fusion) áp dụng cho bài toán phân loại ảnh di sản (từ 2 tập dữ liệu là New Year Print và Clay Figurine). Mô hình này sử dụng cơ chế chú ý đa phương thức (Multimodal Attention) để tập trung vào các vùng ảnh và từ khóa quan trọng, kết hợp với hợp nhất phân cấp (Hierarchical Fusion) để nắm bắt các tương tác động giữa ảnh và văn bản. Kết quả thực nghiệm cho thấy sự vượt trội so với các phương pháp đơn lẻ. Tuy nhiên, giới hạn của nghiên cứu này là chưa khai thác được đặc trưng của âm thanh.\vspace{0.3cm}

Để giải quyết bài toán biểu diễn dữ liệu phức tạp hơn, hướng tiếp cận dựa trên đồ thị đang nhận được sự quan tâm lớn. Trong bài khảo sát toàn diện của Peng và cộng sự \cite{peng2024learningmultimodalgraphssurvey} về học máy trên đồ thị đa phương thức (Learning on Multimodal Graphs), các tác giả khẳng định rằng đồ thị là cấu trúc dữ liệu tự nhiên nhất để biểu diễn các mối quan hệ không đồng nhất. Việc xem mỗi đối tượng dữ liệu (hoặc mỗi phương thức) là một đỉnh trong đồ thị cho phép mô hình nắm bắt được cấu trúc topo và sự phụ thuộc lẫn nhau giữa các nguồn thông tin. Mặc dù vậy, đa số các kiến trúc GNN hiện nay vẫn sử dụng mạng Perceptron đa lớp (MLP) tại các nút để xử lý đặc trưng. Điều này dẫn đến hạn chế về "hộp đen" (khả năng giải thích kém) và hiệu quả tham số chưa tối ưu khi mô hình hóa các hàm phức tạp trong dữ liệu văn hóa.\vspace{0.3cm}

Nhằm khắc phục nhược điểm của MLP, Liu và cộng sự \cite{liu2025kankolmogorovarnoldnetworks} đã giới thiệu một kiến trúc đột phá mang tên Kolmogorov-Arnold Network (KAN). Thay vì sử dụng hàm kích hoạt cố định tại các nút như MLP, KAN đặt các hàm kích hoạt có thể học được trên các cạnh liên kết. Kiến trúc này không chỉ giúp giảm số lượng tham số đáng kể mà còn cung cấp khả năng giải thích minh bạch hơn về mặt toán học. Tuy nhiên, tiềm năng của việc kết hợp KAN vào các mô hình đồ thị đa phương thức (Multimodal Graph + KAN) để xử lý dữ liệu video văn hóa quy mô lớn vẫn là một hướng đi mới mẻ và chưa có nhiều công bố khoa học.\vspace{0.3cm}

Tại Việt Nam, các nghiên cứu nền tảng như PhoBERT của Nguyễn và cộng sự \cite{nguyen-tuan-nguyen-2020-phobert} đã giải quyết tốt bài toán xử lý ngôn ngữ tự nhiên tiếng Việt. Song, các hệ thống tích hợp đa phương thức chuyên sâu cho ICH vẫn chưa được khai thác nhiều. Đa phần các nghiên cứu hiện tại chỉ dừng lại ở việc xử lý riêng lẻ hình ảnh hoặc văn bản, chưa có một phương pháp, kiến trúc thống nhất để đồng bộ hóa và học biểu diễn từ cả video, âm thanh và văn bản tiếng Việt.\vspace{0.3cm}

Tóm lại, từ các phân tích trên, có thể thấy khoảng trống nghiên cứu nằm ở hai điểm chính: (1) Sự thiếu hụt các mô hình xử lý video văn hóa kết hợp cả ba luồng thông tin (thay vì chỉ ảnh và text như MICMLF); và (2) Chưa có sự ứng dụng các kiến trúc mới như KAN vào GNN để tối ưu hóa hiệu năng và khả năng giải thích. Đây chính là động lực để tôi đề xuất mô hình đồ thị đa phương thức kết hợp GNN (GAT) và KAN.